\section{Stability of the original weighted fibers}
\label{sec:v88-fibers}
The preceding estimate controls both signs of physical angular-density
variation. It therefore gives more than a bound on inverse-twist loss:
the original weighted source has a finite-count total-variation modulus
on its complete common-survivor family.

For $0<h\le d(a,\delta)$ let $\mathcal K_{t,h}$ be the members of
$\mathcal E_t(a,\delta)$ with $t+h<b_j$. On this same disjoint
parameter space define two measures
\[
 d\sigma_0=W^\varepsilon(\Psi_j(\theta,t))A_j(t,\theta)\,d\theta,
 \qquad
 d\sigma_h=W^\varepsilon(\Psi_j(\theta,t+h))A_j(t+h,\theta)\,d\theta.
\]
The second measure is the target fiber pulled back by the actual
segment correspondence. It is not the raw target restricted by new
labels, nor an independently renormalized source. We use the full
variation norm, so for these measures it is the integral of the
absolute difference of their densities, with no factor $1/2$.

\begin{theorem}[Original weighted roof-fiber stability]
\label{thm:v88-fiber-stability}
For almost every $(t,h)$ at finite count,
\begin{equation}\label{eq:v88-fiber-stability}
 \|\sigma_h-\sigma_0\|_{\rm var}
 \le C\mathcal V_h^{a,\delta}(t)
                 +K_\Theta\varepsilon^{-1}\sqrt h\,S_t(a,\delta).
\end{equation}
For the thresholds of Theorem~\ref{thm:v88-expanded},
\begin{equation}\label{eq:v88-fiber-finite}
 \sup_{m,R,n,k}m^2\esssup_t\|\sigma_h-\sigma_0\|_{\rm var}
 \le K_{B_\dagger,\Theta}
                   (h^{1/10}+\varepsilon^{-1}h^{1/4}).
\end{equation}
This holds for every fixed $\varepsilon>0$ and all sufficiently small
$h$ simultaneously. It controls every measurable test of absolute
value at most one under the same correspondence. Exited states are
not included in either measure and are charged separately below.
\end{theorem}
\begin{proof}
For each common survivor write $W_s=W^\varepsilon(\Psi(t+s))$.
Since $0\le W_s\le1$,
\[
 |W_hA_h-W_0A_0|\le|A_h-A_0|+A_h|W_h-W_0|.
\]
The full derivative estimate $|A'|\le CAI_m/z$ pays the first
term by $C\mathcal V_h^{a,\delta}$. It includes the change of
$\chi$ as well as $\omega$. The unchanged all-contact guard estimate
in Supplement Corollary~\ref{supp-v78archive:cor:v78-guard} gives
$|W_h-W_0|\le K_\Theta\varepsilon^{-1}\sqrt h$.
Summing $A_h$ over the common-survivor family costs at most $S_t$.
The identities use the common joint coarea kernels; the physical
Jacobian has been incorporated in $A_h$, not set equal to one.
This proves \eqref{eq:v88-fiber-stability}.

Lemma~\ref{lem:v88-trace} and \eqref{eq:v88-uniform-energy} give
$m^2S_t\le Kd^{-1}=Kc_d^{-1}h^{-1/4}$. Combine this with
\eqref{eq:v88-expanded}; multiplication by $\sqrt h$ leaves
$h^{1/4}$. Duality for signed measures proves the assertion for all
bounded measurable tests without selecting different null sets for
each test: equality of the two kernels is already an equality of
measures outside the common null set.
\end{proof}

\subsection{The complete original exceptional source}
Define, on the original fiber and with its original weight,
\begin{equation}\label{eq:v88-residual}
 R_h^{\varepsilon,*}(t)=\int W^\varepsilon
     \one_{\{z<h^{1/8}\ \mathrm{or}\ \ell<h^{1/4}\}}\,d\mu_t,
 \qquad \ell=\tau^-+\tau^+.
\end{equation}
Both alternatives include every endpoint type. The union is counted
once. Write $\mathfrak R_{\varepsilon,M}^*
=\lim_{h\downarrow0}\mathcal H_M(R_h^{\varepsilon,*})$.
Cofinality of the threshold pairs gives exactly
$\mathfrak R_{\varepsilon,M}^*=\mathfrak R_{\varepsilon,M}$
from Supplement Theorem~\ref{supp-v87archive:thm:v87-endpoint}.
Thus the larger controlled family does not by itself change or annul
the ultimate charge.

\begin{proposition}[A finite-count mass bound for the exceptional union]
\label{prop:v88-exceptional-mass}
On any fixed auxiliary band, with $\rho_*<1$ and fixed constants,
\begin{align}
 m^2\int_JR_h^{\varepsilon,*}(t)\,dt
 \le{}&C(h^{\alpha/8}+h^{1/128})(|J|+B_*^{-1})\notag\\
 &+C_{B_*}(1+B_*|J|)m^3\rho_*^{m/2},\qquad \alpha=129/128.
                                       \label{eq:v88-exceptional-mass}
\end{align}
Consequently its fixed-window collision-normalized mass is
$O(|J|h^{1/128})$, uniformly in $\varepsilon$. This assertion is a
mass bound, not a pointwise-height bound.
\end{proposition}
\begin{proof}
Drop $W^\varepsilon\le1$ only on the positive upper comparison.
The initial-gradient part uses the endpoint-only finite-count estimate
in Supplement equation~\eqref{supp-v86archive:eq:v86-endpoint-input}
at $\zeta=h^{1/8}$. If $\ell<h^{1/4}$, the forward lifetime is
less than $h^{1/4}$, regardless of the incoming or outgoing endpoint
type. The complete true-exit starting-mass bound of Supplement
Theorem~\ref{supp-v83archive:thm:v83-exit-local-mass} therefore applies
with width $h^{1/4}$ and contributes $h^{1/128}$. Add the two upper
bounds, which may overcount their intersection only on the upper
side. Enlarge $\rho_*<1$ to combine their remainders. The stated
fixed-window consequence takes the collision limit and then enlarges
the auxiliary band. No claim about the height of the union follows
from this operation.
\end{proof}

\begin{theorem}[The original pointwise budget with an expanded controlled family]
\label{thm:v88-endpoint}
For sufficiently small fixed $h>0$ and $\varepsilon>0$,
\begin{align}
 \mathcal H_M(b^\varepsilon)\le{}&
 \mathcal H_M(R_h^{\varepsilon,*})+C\varepsilon^{1/16}
       +Ch^{1/32}+Kh^{1/10}\notag\\
       &+C_\Theta(1+\varepsilon^{-1})\sqrt h.
                                              \label{eq:v88-endpoint}
\end{align}
With $\varepsilon(B)=A_0B^{-1/12}$, it follows that
\begin{equation}\label{eq:v88-squeeze}
 \mathfrak R_{\varepsilon(B),M}^*\le\mathcal E_M+C_MB^{-1/192},
 \qquad
 \mathcal E_M\le\mathfrak R_{\varepsilon(B),M}^*+C_MB^{-1/192}.
\end{equation}
The remaining assertion for the original scalar theorem is
$\lim_{B\to\infty}\mathfrak R_{\varepsilon(B),M}^*=0$ for every $M$.
It is not proved by the mass or common-survivor stability theorem.
\end{theorem}
\begin{proof}
Split the initial fiber before estimating any exit. On the exceptional
union, survivor and exit are disjoint; their positive defect is at most
the original weight, hence at most $R_h^{\varepsilon,*}$ in total.
On its complement, optical variation pays survivors and genuine grazing
exits by \eqref{eq:v88-expanded}. For a nongrazing first exit,
$\omega(t)\le\omega(b-)+\int_t^b[-\omega']_+$.
The endpoint term, with its physical angular normalization, is bounded
by a fixed multiple of the same outgoing first-exit assignment on
segments of total residence at least $\delta(h)=h^{1/4}$.
Its ordered height is at most $Ch^{1/32}$ by the retained positive
residence theorem, with this positive cutoff fixed before the collision
limit. The pullback Jacobian cost is paid by the variation term.

Insert these bounds into the complete source-shift inequality
\eqref{eq:v83-shift-input}. The original fixed-width bad source and
unweighted source estimates pay its average and guard terms. This proves
\eqref{eq:v88-endpoint}. Here the sharper ordered guard estimate is
used; \eqref{eq:v88-fiber-finite} is the stronger finite-count statement
for arbitrary measurable tests on the common-survivor family.

Now let $h$ decrease after $\mathcal H_M$, keeping $\varepsilon$
fixed. Since $R_h^{\varepsilon,*}\le b^\varepsilon$, the resulting
height differs from its residual charge by at most
$C\varepsilon^{1/16}$. The unchanged positive scalar remainder gives
\eqref{eq:v88-squeeze} by the triangle inequality. All arithmetic
classes remain unnormalized. The prior sufficient charge pairs remain
in the complete archive; this direct-source equivalence neither proves
those larger charges zero nor promotes downstream laws before scalar
closure.
\end{proof}
